\documentclass[10pt,letterpaper]{article}
\usepackage[margin=0.82in]{geometry}
\usepackage{amsmath,amssymb,amsthm,booktabs,array}
\usepackage[T1]{fontenc}
\usepackage{lmodern}
\usepackage{xcolor}
\usepackage{fancyhdr}
\usepackage[colorlinks=true,linkcolor=blue!45!black,urlcolor=blue!45!black]{hyperref}
\setlength{\parindent}{0pt}
\setlength{\parskip}{6pt}
\pagestyle{fancy}
\fancyhf{}
\fancyhead[L]{\small Conjecture 00000001227}
\fancyhead[R]{\small Complete proof in the finite-game setting}
\fancyfoot[C]{\small\thepage}
\setlength{\headheight}{14pt}
\newtheorem{theorem}{Theorem}
\newtheorem{lemma}{Lemma}
\newcommand{\Win}{\mathsf{Win}}
\newcommand{\Ess}{\mathsf{E}}
\newcommand{\code}[1]{\texttt{\small #1}}
\begin{document}

\begin{center}
{\LARGE\bfseries Maximum matchings and vertex geography}\\[7pt]
{\large A Lean 4 proof of conjecture 00000001227}\\[6pt]
C0ldSmi1e \quad $\cdot$ \quad 4 October 2026
\end{center}

\section*{Statement and scope}
For every finite simple undirected graph $G$ and specified initial vertex $v$,
the first player wins normal-play vertex geography precisely when every
maximum-cardinality matching of $G$ saturates $v$. We prove this for all such
graphs, including disconnected graphs and isolated starting vertices.
No perfect matching is assumed to exist.

The accompanying \code{conjecture.md} preserves the exact English and Chinese
statement. Both versions explicitly use \emph{every maximum matching} in the
winning criterion. We interpret a vertex's belonging to a matching as being an
endpoint of a matching edge, and ``maximum'' as maximum edge cardinality,
not merely inclusion-maximal. The introductory mention of perfect matchings
does not replace that quantifier or restrict the class of graphs.

The short source does not explicitly specify finiteness, simplicity, turn order,
or the losing condition. We use the conventional finite, simple, undirected,
impartial normal-play game, as described in the literature
\cite{fox,gates}. These are disclosed named-game conventions, not additional
words attributed to the source. We do not assert an infinite-graph theorem:
infinite play and draws would require a separately specified outcome convention.

\section*{The actual game and matching objects}
The token is already on $v$ before the first move. Players alternate moving it
along an edge to a previously unvisited vertex. The starting vertex counts as
visited. On leaving a vertex it is deleted, so it can never be revisited.
A player with no legal move loses.

A residual position is $(S,v)$, where $S$ is the finite set consisting of the
current vertex and all still-unvisited vertices, and $v\in S$. Its legal moves
are to $(S\setminus\{v\},w)$ with $w\in S\setminus\{v\}$ and $vw\in E(G)$.
Each move decreases $|S|$ by one. Thus ordinary finite backward induction gives
the winning predicate
\[
 \Win(S,v)\ \Longleftrightarrow\
 \exists w\in S\setminus\{v\},\quad vw\in E(G)
 \ \text{and}\ \neg\Win(S\setminus\{v\},w).
\]
This recurrence describes a winning strategy: choose a losing successor on
one's own turn; every opponent move from a losing position returns a winning
position. The decreasing vertex count ensures termination. With no successor,
the existential statement is false, as normal play requires.

A matching on $S$ is a finite set of unordered edges of $G$, all endpoints in
$S$, with distinct edges sharing no endpoint. It saturates $v$ if some edge
contains $v$. It is maximum if its number of edges is at least that of every
matching on $S$. Write $\Ess(S,v)$ when every maximum matching on $S$ saturates
$v$. A maximum always exists: all matchings lie in the finite powerset of the
unordered pairs on $S$, and the empty matching is a candidate.

\newpage
\section*{Proof}
\begin{lemma}[The matching recurrence]
For every valid residual position $(S,v)$,
\[
 \Ess(S,v)\ \Longleftrightarrow\
 \exists w\in S\setminus\{v\},\quad vw\in E(G)
 \ \text{and}\ \neg\Ess(S\setminus\{v\},w).
\]
\end{lemma}
\begin{proof}
Suppose first that $\Ess(S,v)$ holds. Choose a maximum matching $M$ on $S$.
It has an edge $e=vw$ incident to $v$. Since the graph has no loops,
$w\in S\setminus\{v\}$. Put $N=M\setminus\{e\}$.
The matching property ensures that $N$ saturates neither $v$ nor $w$;
in particular, it is a matching on $S\setminus\{v\}$.

We claim that $N$ is maximum there. Any matching $L$ on $S\setminus\{v\}$
is also a matching on $S$, so $|L|\leq |M|$. Equality would make $L$ a maximum
matching on $S$ that does not saturate $v$, contradicting $\Ess(S,v)$.
Therefore $|L|<|M|$, and the integer inequality implies
$|L|\leq |M|-1=|N|$. Thus $N$ is maximum on $S\setminus\{v\}$ and misses $w$,
which proves $\neg\Ess(S\setminus\{v\},w)$.

Conversely, suppose that $vw$ is a legal move and
$\neg\Ess(S\setminus\{v\},w)$. There is a maximum matching $N$ on
$S\setminus\{v\}$ that misses $w$. It also misses $v$, so
$K=N\cup\{vw\}$ is a matching on $S$ with $|K|=|N|+1$.

Let $M$ be any maximum matching on $S$. If it missed $v$, it would itself be
a matching on $S\setminus\{v\}$, giving $|M|\leq |N|$ by maximal cardinality
of $N$. But maximum cardinality of $M$ on $S$ gives
$|N|+1=|K|\leq |M|$, a contradiction. Every such $M$ must saturate $v$.
\end{proof}

\begin{theorem}[The vertex-geography criterion]
At every valid residual position $(S,v)$,
$\Win(S,v)\longleftrightarrow\Ess(S,v)$.
Consequently the first player wins from any specified vertex of any finite
simple undirected graph if and only if every maximum matching saturates it.
\end{theorem}
\begin{proof}
Induct on the size of $S$. The game recurrence and the matching recurrence
have identical legal successors. Each successor has smaller available-vertex
set and a valid current vertex. The induction hypothesis identifies the two
predicates at each successor, and hence identifies their negations and the
existential statements defining the current position. This establishes the
equivalence. Initially take $S$ to be the entire vertex set.
\end{proof}

An isolated start has no legal successor and therefore loses, regardless of
other components. The Lean project also states this case separately. When the
vertex type is empty there is no start vertex; the universal initial theorem
remains valid without needing a nonempty-graph assumption.

\section*{Attribution}
The characterization is already a known theorem. Gates et al.,
Theorem 2.8, explicitly attribute it to Theorem 1.1 of
Fraenkel, Scheinerman and Ullman's \emph{Undirected edge geography}
\cite{gates,fsu}; Fox and Geissler also state the vertex criterion
\cite{fox}. The source's description as an extension of an edge result is
historically misleading. This submission gives a complete proof and formalization
of its mathematical criterion; it makes no novelty claim and does not treat
that historical sentence as a further Lean theorem.

\newpage
\section*{Formalization and verification}
The project pins Lean 4.19.0 and Mathlib v4.19.0. Graphs are Mathlib
\code{SimpleGraph} objects. Edges are \code{Sym2} unordered pairs;
matchings are actual finite edge sets, and size counts edges.
The seven authored definitions and fourteen theorems are distributed as follows.

\begin{center}\small
\begin{tabular}{@{}p{0.27\linewidth}p{0.67\linewidth}@{}}
\toprule
Module & Mathematical content \\
\midrule
\code{Matching.lean} & Matching, saturation, maximum cardinality, the universal
matching condition, existence of a maximum, and the deletion recurrence. \\
\code{Game.lean} & Legal moves and a separate well-founded definition of
\code{Wins}, its normal-play recurrence and uniqueness, and the initial position. \\
\code{Main.lean} & Equivalence at every valid residual position, the full initial
criterion, and loss from an isolated start. \\
\bottomrule
\end{tabular}
\end{center}

The main declaration, in namespace \code{Conjecture1227}, is
\begin{center}\small\ttfamily
first\_player\_wins\_iff\_every\_maximum\_matching\_saturates.
\end{center}
It quantifies over an arbitrary finite vertex type, arbitrary simple graph,
arbitrary start and every maximum matching. The \code{Fintype} and
\code{DecidableEq} instances provide finite-set representation, with no graph
restriction. The stronger \code{wins\_iff\_essential} theorem allows an arbitrary
finite residual set even in an otherwise unrestricted ambient vertex type.
Invalid positions $v\notin S$ are assigned false by the game definition, but
the criterion is asserted only for valid positions. Initially and after every
legal move the validity condition holds.

\code{Wins} is defined only through legal moves and normal play; it has no
matching input. The general uniqueness theorem for this terminating recurrence
is applied to the independently proved matching recurrence. The proof uses
neither a matching oracle nor a hypothesis assuming the desired equivalence.

From the bundled \code{lean/} directory, run \code{lake build}, then
\code{lake env lean -DwarningAsError=true Check.lean}. The library build itself
treats warnings as errors. \code{Check.lean} prints all authored definitions
and every authored theorem's type and transitive axiom dependencies.
The accompanying verification records document an independent build with fresh
local outputs, strict replays of every source file, all nine pinned dependency
revisions, compiler identity, source hashes, and a complete inventory of constants
originating in the submitted modules, including generated declarations.
Only the standard foundational axioms are permitted. No admitted proof,
\code{native\_decide}, custom axiom or authored unsafe proof is used.

As auxiliary evidence, run
\code{python3 auxiliary/check\_vertex\_geography.py -{}-max-n 6} from the submission
directory. Separate algorithms enumerate actual matchings and evaluate game
trees. All 33,868 labelled simple graphs on zero through six vertices and
202,013 starts agree, including 6,391 disconnected graphs and 6,505 isolated
starts. The empty graph contributes no start. This exhaustive finite check is
supplementary; the theorem above and its Lean proof cover every finite size.

The package includes the exact bilingual source, the report and matching PDF,
the Lean project, the auxiliary script and results, and verification records.
An independent internal semantic review checks their correspondence.
Local verification and internal review are distinct from repository maintainer
acceptance.

\begin{thebibliography}{9}\small
\bibitem{fox} N. Fox and C. Geissler.
\emph{On the Computational Complexities of Various Geography Variants}, 2021,
\S1 and \S5. \url{https://arxiv.org/abs/2108.09367}.
\bibitem{gates} Z. Gates, D. Hayes, R. Kelvey, W. Lillis and M. Phan.
\emph{Cycle Games on Graphs}. Integers 24 (2024), article G4, \S2.1, Theorem 2.8.
\url{https://math.colgate.edu/~integers/yg4/yg4.pdf}.
\bibitem{fsu} A. S. Fraenkel, E. R. Scheinerman and D. Ullman.
\emph{Undirected edge geography}. Theoretical Computer Science 112 (1993),
371--381. \url{https://doi.org/10.1016/0304-3975(93)90026-P}.
The vertex-theorem attribution is corroborated by \cite{gates}; the original
full text was not inspected for this report.
\end{thebibliography}
\end{document}
